% ===========================================================================
% Formal Semantics of the NeuroPLC Intermediate Representation
% ===========================================================================
\subsection{Formal Semantics of the NeuroPLC IR}
\label{sec:ir-semantics}

The NeuroPLC IR is not merely a compiler data structure---it is a
\textbf{formal language} with a defined syntax, operational semantics,
and type system. This formalization serves two purposes:
(1)~it makes Theorem~\ref{thm:compiler} (compiler correctness)
precise as a \textbf{type soundness} theorem in the style of programming
language theory~\cite{milner1978theory,pierce2002types}; and
(2)~it establishes that the SVNN conditions
(Definition~\ref{def:svnn}, Theorem~\ref{thm:svnn_conditions}) constitute
a \textbf{type system} for neural computation graphs.

\subsubsection{Syntax}

The IR language $\mathcal{L}_{\text{IR}}$ is defined by the following BNF grammar:

\begin{equation}
\begin{aligned}
e \;::=&\; \texttt{Input}(x)                                    &&\text{input variable} \\
\mid&\; \texttt{MatMul}(W,\; e)                                  &&\text{matrix multiplication} \\
\mid&\; \texttt{BsplineLUT}(T,\; g,\; e)                         &&\text{B-spline LUT evaluation} \\
\mid&\; \texttt{StandardAct}(\sigma,\; e)                        &&\text{standard activation (SiLU)} \\
\mid&\; \texttt{Add}(e_1,\; e_2)                                 &&\text{element-wise addition} \\
\mid&\; \texttt{Softmax}(e)                                      &&\text{softmax normalization} \\
\mid&\; \texttt{Argmax}(e)                                       &&\text{argmax classification}
\end{aligned}
\label{eq:ir_bnf}
\end{equation}

where $W \in \mathbb{R}^{m \times n}$ is a weight matrix,
$T \in \mathbb{R}^{N}$ is a precomputed LUT array,
$g \in \mathbb{R}^{N_{\text{grid}}}$ is a knot grid,
$\sigma \in \{\text{SiLU}\}$, and $x$ denotes a named input variable.
The grammar is \textit{closed under the SVNN operation set}: every
production corresponds to either a pure linear operation (MatMul, Add)
or a pure element-wise operation (BsplineLUT, StandardAct, Softmax,
Argmax), satisfying Condition~1 by construction.

\subsubsection{Denotational Semantics}

We define two interpretation functions for each IR expression:

\begin{itemize}
    \item $\llbracket e \rrbracket_{\mathbb{R}}$: the
    \textbf{real-valued semantics}---the mathematical function computed
    by $e$ under exact real arithmetic (matching PyTorch's float32
    training semantics up to IEEE~754 roundoff).

    \item $\llbracket e \rrbracket_{\text{SCL}}$: the
    \textbf{SCL execution semantics}---the function computed by the
    compiled SCL code on a Siemens PLC under IEC~61131-3 REAL
    arithmetic (IEEE~754 binary32 with soft-float execution).
\end{itemize}

The real-valued semantics $\llbracket \cdot \rrbracket_{\mathbb{R}}$
is defined by structural induction on $e$:

\begin{align}
\llbracket \texttt{Input}(x) \rrbracket_{\mathbb{R}}(\rho) &= \rho(x) \\
\llbracket \texttt{MatMul}(W, e) \rrbracket_{\mathbb{R}}(\rho) &= W \cdot \llbracket e \rrbracket_{\mathbb{R}}(\rho) \\
\llbracket \texttt{BsplineLUT}(T, g, e) \rrbracket_{\mathbb{R}}(\rho) &=
\text{interp}\bigl(T, g, \text{clamp}(\llbracket e \rrbracket_{\mathbb{R}}(\rho),
[g_0, g_{N_g-1}])\bigr) \\
\llbracket \texttt{StandardAct}(\text{SiLU}, e) \rrbracket_{\mathbb{R}}(\rho) &=
\text{SiLU}(\llbracket e \rrbracket_{\mathbb{R}}(\rho)) \\
\llbracket \texttt{Add}(e_1, e_2) \rrbracket_{\mathbb{R}}(\rho) &=
\llbracket e_1 \rrbracket_{\mathbb{R}}(\rho) + \llbracket e_2 \rrbracket_{\mathbb{R}}(\rho) \\
\llbracket \texttt{Softmax}(e) \rrbracket_{\mathbb{R}}(\rho) &=
\text{softmax}(\llbracket e \rrbracket_{\mathbb{R}}(\rho)) \\
\llbracket \texttt{Argmax}(e) \rrbracket_{\mathbb{R}}(\rho) &=
\argmax_i\; \llbracket e \rrbracket_{\mathbb{R}}(\rho)_i
\end{align}

where $\rho: \text{Var} \to \mathbb{R}^{d_{\text{in}}}$ is an
environment mapping input variables to their real-valued vectors,
$\text{interp}(T, g, x)$ performs piecewise-linear interpolation of
table $T$ on grid $g$ at point $x$, and $\text{clamp}(x, [a,b])$
clamps $x$ to the validated domain $[a,b]$.

\subsubsection{Type System}

The NeuroPLC IR is equipped with a simple type system whose judgments
have the form $\Gamma \vdash e : \tau$, where types
$\tau \in \{\texttt{linear}, \texttt{elemwise}, \texttt{svnn}\}$
classify IR nodes by their operation category.

\begin{center}
\framebox{
\begin{minipage}{0.80\linewidth}
\vspace{4pt}
\begin{prooftree}
\AxiomC{}
\UnaryInfC{$\Gamma \vdash \texttt{Input}(x) : \texttt{linear}$}
\DisplayProof
\quad\textsc{[T-Input]}
\end{prooftree}

\vspace{6pt}
\begin{prooftree}
\AxiomC{$\Gamma \vdash e : \texttt{linear}$}
\UnaryInfC{$\Gamma \vdash \texttt{MatMul}(W, e) : \texttt{linear}$}
\DisplayProof
\quad\textsc{[T-MatMul]}
\end{prooftree}

\vspace{6pt}
\begin{prooftree}
\AxiomC{$\Gamma \vdash e : \texttt{linear}$}
\AxiomC{$M_2(T,g) < \infty$}
\BinaryInfC{$\Gamma \vdash \texttt{BsplineLUT}(T, g, e) : \texttt{elemwise}$}
\DisplayProof
\quad\textsc{[T-BsplineLUT]}
\end{prooftree}

\vspace{6pt}
\begin{prooftree}
\AxiomC{$\Gamma \vdash e_1 : \texttt{linear}$}
\AxiomC{$\Gamma \vdash e_2 : \texttt{elemwise}$}
\BinaryInfC{$\Gamma \vdash \texttt{Add}(e_1, e_2) : \texttt{svnn}$}
\DisplayProof
\quad\textsc{[T-Add]}
\end{prooftree}

\vspace{6pt}
\begin{prooftree}
\AxiomC{$\Gamma \vdash e : \texttt{svnn}$}
\UnaryInfC{$\Gamma \vdash \texttt{Softmax}(e) : \texttt{svnn}$}
\DisplayProof
\quad\textsc{[T-Softmax]}
\end{prooftree}
\vspace{4pt}
\end{minipage}
}
\end{center}

The typing rules directly encode SVNN Conditions~1--2:
\textsc{[T-BsplineLUT]} requires $M_2(T,g) < \infty$ (Condition~2:
computable curvature bound); \textsc{[T-Add]} requires one argument
of linear type and one of element-wise type, enforcing the dual-path
structure of KAN layers; \textsc{[T-MatMul]} preserves linear type,
reflecting Condition~1 (linear operations remain linear); and
\textsc{[T-Softmax]} propagates the SVNN type to the output.

\vspace{4pt}
\begin{kingthm}[IR Type Soundness]
\label{thm:ir-soundness}
If $\vdash e : \texttt{svnn}$ (i.e., $e$ is well-typed under the
IR type system), then the compiled SCL program $C(e)$ satisfies:
\begin{equation}
\forall \rho \in \text{Domain}(X):\;
\|\llbracket e \rrbracket_{\mathbb{R}}(\rho) -
\llbracket C(e) \rrbracket_{\text{SCL}}(\rho)\|_\infty
\leq \varepsilon(e)
\label{eq:type_soundness}
\end{equation}
where $\varepsilon(e)$ is the DA-computed error bound
(Theorem~\ref{thm:svnn_conditions}, Eq.~\ref{eq:total_bound}).
\end{kingthm}

\vspace{2pt}
\noindent\textit{Proof.}
By structural induction on the typing derivation of $e$.
Each typing rule corresponds to a proof obligation in
Theorem~\ref{thm:compiler}:

\begin{itemize}
    \item \textsc{[T-Input]}: Input carries no error by construction
    ($\Delta^{(0)} = 0$).

    \item \textsc{[T-MatMul]}: Linear layer error amplification is
    bounded by $\|W\|_{1,\infty} \cdot \|\delta\|_\infty$
    (Theorem~\ref{thm:svnn_conditions}, Step~1). DA improves the
    constant via sign-structural tightening.

    \item \textsc{[T-BsplineLUT]}: The $M_2 < \infty$ premise
    guarantees the de~Boor LUT error bound
    $M_2 \cdot h^2 / 8$ (Eq.~\ref{eq:de_boor_general}).

    \item \textsc{[T-Add]}: Addition introduces no new error
    sources; the error is the sum of branch errors.

    \item \textsc{[T-Softmax]}: Softmax is a contraction mapping
    in $\ell_\infty$ (Lipschitz constant $\leq 1$)---it does not
    amplify upstream errors.
\end{itemize}

The induction composes these per-node bounds into the
end-to-end guarantee of Eq.~\ref{eq:type_soundness}, exactly
recovering Theorem~\ref{thm:compiler} but now
phrased as a \textbf{type soundness} result: well-typed IR
programs compile to SCL with bounded semantic deviation. $\square$

\vspace{4pt}
\noindent\textbf{Relationship to Programming Language Theory.}
Theorem~\ref{thm:ir-soundness} is the neural-network analogue of the
canonical type soundness theorem in programming languages~\cite{milner1978theory}:
``well-typed programs cannot go wrong.'' In our setting, ``going wrong''
means producing an output whose deviation from the trained model exceeds
the certified bound $\varepsilon$. The SVNN type system guarantees that
this cannot happen---a property that standard MLP architectures lack
because they are \textit{ill-typed} under these rules (the coupled
$\sigma(W\mathbf{x} + \mathbf{b})$ operation has no typing rule).

This perspective reframes the central contribution of NeuroPLC: it is
not merely a compiler, but a \textbf{type-theoretic framework} for
certifiable neural architectures, in which the PLC deployment is one
concrete instantiation of a more general principle---that
\textit{well-typedness of an architecture implies certifiability of
its compilation}.
